% Generated draft: requires mathematical and source review.
\documentclass[11pt]{article}
\usepackage[margin=25mm]{geometry}
\usepackage{fontspec}
\setmainfont{Latin Modern Roman}
\usepackage{amsmath,amssymb,mathtools,unicode-math}
\setmathfont{Latin Modern Math}
\usepackage{enumitem,needspace,xurl,hyperref}
\setlength{\parindent}{0pt}
\setlength{\parskip}{5pt}
\setlength{\emergencystretch}{3em}
\newcommand{\N}{\mathbb{N}}
\newcommand{\Z}{\mathbb{Z}}
\newcommand{\Q}{\mathbb{Q}}
\newcommand{\R}{\mathbb{R}}
\newcommand{\C}{\mathbb{C}}
\newcommand{\F}{\mathbb{F}}
\newcommand{\A}{\mathbb{A}}
\newcommand{\PP}{\mathbb{P}}
\newcommand{\E}{\mathbb{E}}
\newcommand{\eps}{\varepsilon}
\newcommand{\dd}{\,\mathrm{d}}
\begin{document}
\section*{UnsolvedMath ID: 700115}
\textbf{Source problem code:} catalogue entry
\subsection{Definitions and mathematical statement}
For a fixed finitary algebraic signature \(L\), a quasi-identity is a universally quantified implication
\[
\forall x_1\cdots\forall x_n\;\left(\bigwedge_{i=1}^{m}s_i(\mathbf{x})=t_i(\mathbf{x})\;\Longrightarrow\;s(\mathbf{x})=t(\mathbf{x})\right),
\]
where all displayed expressions are \(L\)-terms and \(m\geq 0\). The empty conjunction is interpreted as true. In the usual universal-algebraic terminology, a quasivariety is a class of algebras axiomatized by a set of quasi-identities. Once an ambient class \(\mathcal U\) and its signature have been specified, the corresponding relative notion is
\[
\mathcal Q_L(\mathcal U)=\left\{\{A\in\mathcal U:A\models\Sigma\}:\Sigma\text{ is a set of }L\text{-quasi-identities}\right\}.
\]
This notation supplies a possible object collection; it does not supply a groupoid structure.

\textbf{Source question (A. I. Mal'cev, 1967).} Determine the structure of the groupoid of quasivarieties in each of the following cases:
\begin{enumerate}
\item[(a)] quasivarieties of semigroups;
\item[(b)] quasivarieties of rings;
\item[(c)] quasivarieties of associative rings.
\end{enumerate}
A semigroup has an associative binary multiplication. The source does not specify the ring signatures or conventions, including whether rings have a multiplicative identity and, if so, whether that identity is a named constant. The separate listing of rings and associative rings suggests that part (b) may permit nonassociative multiplication, but this interpretation is not explicitly supplied.

\textbf{Unresolved qualification.} The supplied record does not define its use of \emph{groupoid}. In the binary-operation sense, the operation on quasivarieties is missing; in the categorical sense, the morphisms and composition are missing. The record also does not explicitly resolve the absolute versus relative meaning of quasivariety or conventions concerning empty algebras. Consequently, the three source questions can be preserved, but a uniquely specified mathematical structure to be determined cannot be reconstructed from the supplied evidence alone.
\subsection{Short English statement}
Mal'cev asks for the structure of a groupoid formed from quasivarieties of semigroups, rings, and associative rings. The supplied text leaves the groupoid structure and some algebraic conventions undefined, so the intended classification target requires clarification from the cited sources. The registry labels the problem open; that label has not been independently verified.
\subsection{Sources}
\begin{itemize}
\item Supplied local UnsolvedMath registry record: permanent ID \(700115\), exact source code \texttt{catalogue entry}, title \emph{Kourovka Notebook Problem 2.42}. The supplied status is \texttt{open}; this is a registry label, not an independently verified assessment.
\item \emph{The Kourovka Notebook (Unsolved Problems in Group Theory)}, problem 2.42, attributed to A. I. Mal'cev, second issue (1967). The supplied background also identifies the twenty-first issue edition (2026), \texttt{arXiv:1401.0300v48}. The supplied collection description names E. I. Khukhro and V. D. Mazurov as editors and describes coverage of Issues 1--21 (1965--2026).
\item The statement supplies the comparison reference: A. I. Mal'cev, \emph{Siberian Math. J.}, volume \(8\), no. \(2\) (1967), pp.\ 254--267. No article title is supplied. This reference is presented as a comparison, without an explicit definition of the groupoid in the supplied record.
\item These bibliographic references are supplied attribution only; their contents and the intended formulation have not been independently checked.
\end{itemize}
\end{document}
